\section{Bab 2: Bukti Beriman}
\subsection{Memenuhi Janji, Mensyukuri Nikmat, Menjaga Lisan, dan Menutup Aib}
\[
\text{Iman}
\begin{cases}
  \text{Qalbu (Hati)} \\
  \text{Qoul (Lisan)} \\
  \text{Af'al (Perbuatan)}
\end{cases}
\]

\[
\text{Syu'abul Iman (Cabang-cabang iman)}
\begin{cases}
  \text{Memenuhi janji} \\
  \text{Mensyukuri nikmat} \\
  \text{Menjaga lisan} \\
  \text{Menutup aib}
\end{cases}
\]

Janji adalah kesanggupan seseorang untuk melakukan/tidak melakukan sesuatu yang disampaikan kepada orang lain.


Memenuhi janji berarti melaksanakan sesuatu sesuai dengan kesanggupan yang telah diucapkan.

\subsection{Dalil Memenuhi Janji}
\begin{itemize}
\item \textit{Wa lā taqrabū mālal-yatīmi illā billatī hiya aḥsanu ḥattā yabluġa asyuddah(ū), wa aufū bil-‘ahdi, innal-‘ahda kāna mas'ūlā.} \\
  \textbf{Artinya: }Dan janganlah kamu mendekati harta anak yatim, kecuali dengan cara yang lebih bermanfaat sampai dia mencapai usia dewasa. Dan penuhilah janji, karena janji itu pasti diminta pertanggungjawabannya. \textbf{Al-Isra (34)}

\item \textit{Yā ayyuhal-lażīna āmanū aufū bil-‘uqūd(i), uḥillat lakum bahīmatul-an‘āmi illā mā yutlā ‘alaikum gaira muḥilliṣ-ṣaidi wa antum ḥurum(un), innallāha yaḥkumu mā yurīd.} \\
  \textbf{Artinya: }Wahai orang-orang yang beriman! Penuhilah janji-janji. Hewan ternak dihalalkan bagimu, kecuali yang akan disebutkan kepadamu, dengan tidak manghalalkan berburu ketika kamu sedang berihram (haji atau umrah). Sesungguhnya Allah menetapkan hukum sesuai yang ia kehendaki \textbf{Al-Maidah (1)}
\end{itemize}

\subsubsection{Ciri-Ciri Orang yang Menepati Janji}
\begin{enumerate}
\item Tidak mudah mengucapkan janji.
\item Berpikir sebelum berjanji.
\item Berusaha melaksanakan janji.
\item Bertanggung jawab apabila terjadi kendala.
\item Meminta maaf jika tidak mampu memenuhi janji.
\item Tidak mencari cari alasan.
\end{enumerate}

\subsubsection{Hikmah Menepati Janji}
\begin{enumerate}
\item Mendapat kepercayaan.
\item Membuka tanggung jawab.
\item Membentuk pribadi amanah.
\item Mempererat persaudaraan.
\item Menciptakan kehidupan sosial yang harmonis.
\end{enumerate}

\subsubsection{Bahaya Mengingkari Janji}
\begin{enumerate}
\item Menghilangkan kepercayaan.
\item Merusak hubungan.
\item Menimbulkan konflik.
\item Membentuk karakter tidak bertanggung jawab.
\end{enumerate}

\begin{quotation}
  ``Rasulullah SAW bersabda: ``Tanda-tanda orang munafik itu ada 3 jika berbicara ia berdusta, jika berjanji ia mengingkari, dan jika diberi amanat ia berkhianat.'''' \textbf{(HR Bukhari \& Muslim)}
\end{quotation}
